% FILE: CSE-26_CT-2.tex
\documentclass{article}
\usepackage{amsmath}
\usepackage{amssymb}
\begin{document}

%%%%%%%%%%%%%%%%%%%%%%%%%%%%%%%%%%%%%%%%%%%%%%%%%%
% QUESTION_START
%%%%%%%%%%%%%%%%%%%%%%%%%%%%%%%%%%%%%%%%%%%%%%%%%%
% id=cse26-ct2-seca-q1
% discipline=CSE
% year=2026
% exam_type=CT
% section=A
% ct_number=2
% question_number=1
% topic=Maxima & Minima
% subtopic=General
% difficulty=Medium
% length=Long
% question_type=New
% frequency=1
% appearances=
% OCR_STATUS=VERIFIED
\section*{CT2 (Sec A) - Q1}
Question:
Distinguish between extremum points and inflection points of $f(x)=0$. Find the extremum values of $y = e^x + 2\cos x + e^{-x}$ for $0 < x < 2\pi$.

Solution:
Step 1: Distinguish between extremum points and inflection points.
\begin{itemize}
    \item \textbf{Extremum Point}: A point on a curve $y = f(x)$ where the function reaches a local maximum or local minimum. At an extremum point, the first derivative is zero ($f'(x) = 0$, stationary points) or undefined (critical points), and the first derivative $f'(x)$ changes its sign as it passes through the point.
    \item \textbf{Inflection Point}: A point on a curve $y = f(x)$ where the concavity changes (i.e., the curve changes from concave up to concave down, or vice versa). At an inflection point, the second derivative is zero ($f''(x) = 0$) or undefined, and the second derivative $f''(x)$ changes its sign as it passes through the point.
\end{itemize}

Step 2: Find the stationary/critical points of $y = e^x + 2\cos x + e^{-x}$ for $0 < x < 2\pi$.
We rewrite the function as:
\[
y = 2\cosh x + 2\cos x
\]
Differentiating with respect to $x$:
\[
y' = e^x - 2\sin x - e^{-x} = 2\sinh x - 2\sin x
\]
To find the stationary points, set $y' = 0$:
\[
2\sinh x - 2\sin x = 0 \implies \sinh x = \sin x
\]

Step 3: Analyze the equation $\sinh x = \sin x$ on the open interval $(0, 2\pi)$.
For any $x > 0$, we have the inequality:
\[
\sinh x = x + \frac{x^3}{3!} + \frac{x^5}{5!} + \dots > x > \sin x
\]
Since $\sinh x > \sin x$ for all $x > 0$, there are no real solutions to $\sinh x = \sin x$ on the interval $(0, 2\pi)$.

Step 4: Conclude on extremum values.
Since there are no critical points in the open interval $(0, 2\pi)$, the function has no local extremum values on this domain.

Final Answer:
\[
\boxed{\text{There are no extremum values for } y = e^x + 2\cos x + e^{-x} \text{ in the interval } 0 < x < 2\pi.}
\]
%%%%%%%%%%%%%%%%%%%%%%%%%%%%%%%%%%%%%%%%%%%%%%%%%%
% QUESTION_END
%%%%%%%%%%%%%%%%%%%%%%%%%%%%%%%%%%%%%%%%%%%%%%%%%%

%%%%%%%%%%%%%%%%%%%%%%%%%%%%%%%%%%%%%%%%%%%%%%%%%%
% QUESTION_START
%%%%%%%%%%%%%%%%%%%%%%%%%%%%%%%%%%%%%%%%%%%%%%%%%%
% id=cse26-ct2-seca-q2
% discipline=CSE
% year=2026
% exam_type=CT
% section=A
% ct_number=2
% question_number=2
% topic=Applications of Derivatives
% subtopic=General
% difficulty=Medium
% length=Medium
% question_type=New
% frequency=1
% appearances=
% OCR_STATUS=VERIFIED
\section*{CT2 (Sec A) - Q2}
Question:
What is curvature? Find the radius of curvature of the curve $y = \frac{12}{x}$ at $(3, 4)$.

Solution:
Step 1: Define curvature and radius of curvature.
\begin{itemize}
    \item \textbf{Curvature} ($\kappa$): Curvature is a measure of how sharply a curve bends at a given point. For a Cartesian curve $y = f(x)$, it is defined as:
    \[
    \kappa = \frac{|y''|}{\left(1 + (y')^2\right)^{3/2}}
    \]
    \item \textbf{Radius of Curvature} ($\rho$): The radius of curvature is the reciprocal of the curvature:
    \[
    \rho = \frac{1}{\kappa} = \frac{\left(1 + (y')^2\right)^{3/2}}{|y''|}
    \]
\end{itemize}

Step 2: Find the derivatives of the curve $y = \frac{12}{x}$.
\[
y' = \frac{d}{dx}\left(\frac{12}{x}\right) = -\frac{12}{x^2}
\]
\[
y'' = \frac{d^2}{dx^2}\left(\frac{12}{x}\right) = \frac{24}{x^3}
\]

Step 3: Evaluate the derivatives at the point $(3, 4)$.
At $x = 3$:
\[
y' = -\frac{12}{3^2} = -\frac{12}{9} = -\frac{4}{3}
\]
\[
y'' = \frac{24}{3^3} = \frac{24}{27} = \frac{8}{9}
\]

Step 4: Compute the radius of curvature $\rho$.
\[
\rho = \frac{\left[1 + \left(-\frac{4}{3}\right)^2\right]^{3/2}}{\left|\frac{8}{9}\right|} = \frac{\left[1 + \frac{16}{9}\right]^{3/2}}{\frac{8}{9}}
\]
\[
\rho = \frac{\left(\frac{25}{9}\right)^{3/2}}{\frac{8}{9}} = \frac{\frac{125}{27}}{\frac{8}{9}} = \frac{125}{27} \times \frac{9}{8} = \frac{125}{24}
\]

Final Answer:
\[
\boxed{\rho = \frac{125}{24}}
\]
%%%%%%%%%%%%%%%%%%%%%%%%%%%%%%%%%%%%%%%%%%%%%%%%%%
% QUESTION_END
%%%%%%%%%%%%%%%%%%%%%%%%%%%%%%%%%%%%%%%%%%%%%%%%%%

%%%%%%%%%%%%%%%%%%%%%%%%%%%%%%%%%%%%%%%%%%%%%%%%%%
% QUESTION_START
%%%%%%%%%%%%%%%%%%%%%%%%%%%%%%%%%%%%%%%%%%%%%%%%%%
% id=cse26-ct2-seca-q3
% discipline=CSE
% year=2026
% exam_type=CT
% section=A
% ct_number=2
% question_number=3
% topic=Applications of Derivatives
% subtopic=Tangent & Normal
% difficulty=Medium
% length=Medium
% question_type=New
% frequency=1
% appearances=
% OCR_STATUS=VERIFIED
\section*{CT2 (Sec A) - Q3}
Question:
Define sub tangent and sub normal. Find the equation of tangent line and normal line at $(2, 2)$ of the curve $xy^2 = 4(4-x)$.

Solution:
Step 1: Define subtangent and subnormal.
Let $P(x_1, y_1)$ be a point on a curve. Let the tangent and the normal drawn at $P$ intersect the $x$-axis at $T$ and $N$ respectively, and let $M$ be the projection of $P$ on the $x$-axis.
\begin{itemize}
    \item \textbf{Subtangent}: The segment $TM$ on the $x$-axis, whose length is given by:
    \[
    \text{Length of Subtangent} = \left| \frac{y_1}{y'} \right|
    \]
    \item \textbf{Subnormal}: The segment $MN$ on the $x$-axis, whose length is given by:
    \[
    \text{Length of Subnormal} = |y_1 y'|
    \]
\end{itemize}

Step 2: Find the derivative $y'$ of the curve $xy^2 = 4(4-x)$ at $(2, 2)$.
We can rewrite the equation as:
\[
xy^2 = 16 - 4x
\]
Differentiating implicitly with respect to $x$:
\[
y^2 + x \cdot 2y y' = -4
\]
At $(2, 2)$, substitute $x = 2$ and $y = 2$:
\[
(2)^2 + 2(2)(2)y' = -4 \implies 4 + 8y' = -4
\]
\[
8y' = -8 \implies y' = -1
\]
Thus, the slope of the tangent line is $m = -1$.

Step 3: Find the equation of the tangent line at $(2, 2)$.
Using the point-slope form with $m = -1$:
\[
y - 2 = -1(x - 2) \implies y - 2 = -x + 2 \implies x + y - 4 = 0
\]

Step 4: Find the equation of the normal line at $(2, 2)$.
The slope of the normal line is $m' = -\frac{1}{m} = 1$.
Using the point-slope form with $m' = 1$:
\[
y - 2 = 1(x - 2) \implies y - 2 = x - 2 \implies x - y = 0
\]

Final Answer:
\[
\boxed{\text{Tangent: } x + y - 4 = 0, \quad \text{Normal: } x - y = 0}
\]
%%%%%%%%%%%%%%%%%%%%%%%%%%%%%%%%%%%%%%%%%%%%%%%%%%
% QUESTION_END
%%%%%%%%%%%%%%%%%%%%%%%%%%%%%%%%%%%%%%%%%%%%%%%%%%

%%%%%%%%%%%%%%%%%%%%%%%%%%%%%%%%%%%%%%%%%%%%%%%%%%
% QUESTION_START
%%%%%%%%%%%%%%%%%%%%%%%%%%%%%%%%%%%%%%%%%%%%%%%%%%
% id=cse26-ct2-seca-q4
% discipline=CSE
% year=2026
% exam_type=CT
% section=A
% ct_number=2
% question_number=4
% topic=Partial Derivatives
% subtopic=Euler's Theorem
% difficulty=Hard
% length=Long
% question_type=New
% frequency=1
% appearances=
% OCR_STATUS=VERIFIED
\section*{CT2 (Sec A) - Q4}
Question:
State Euler theorem for homogeneous function of two variables with example. If $u = \ln\sqrt{x^2 + y^2 + z^2}$ then show that $(x^2 + y^2 + z^2)\left(\frac{\partial^2 u}{\partial x^2} + \frac{\partial^2 u}{\partial y^2} + \frac{\partial^2 u}{\partial z^2}\right) = 1$.

Solution:
Step 1: State Euler's theorem for homogeneous functions.
\begin{itemize}
    \item \textbf{Euler's Theorem}: If $f(x, y)$ is a homogeneous function of degree $n$ in the independent variables $x$ and $y$, then:
    \[
    x \frac{\partial f}{\partial x} + y \frac{\partial f}{\partial y} = n f(x, y)
    \]
    \item \textbf{Example}: Let $f(x, y) = x^2 + 3xy + y^2$. This is homogeneous of degree $n = 2$.
    \begin{align*}
        \frac{\partial f}{\partial x} = 2x + 3y &\implies x \frac{\partial f}{\partial x} = 2x^2 + 3xy \\
        \frac{\partial f}{\partial y} = 3x + 2y &\implies y \frac{\partial f}{\partial y} = 3xy + 2y^2
    \end{align*}
    Summing them:
    \[
    x \frac{\partial f}{\partial x} + y \frac{\partial f}{\partial y} = 2x^2 + 6xy + 2y^2 = 2(x^2 + 3xy + y^2) = 2 f(x, y)
    \]
    This verifies the theorem.
\end{itemize}

Step 2: Simplify the function $u = \ln\sqrt{x^2 + y^2 + z^2}$.
\[
u = \ln(x^2 + y^2 + z^2)^{1/2} = \frac{1}{2} \ln(x^2 + y^2 + z^2)
\]

Step 3: Calculate the first partial derivatives.
\[
\frac{\partial u}{\partial x} = \frac{1}{2} \cdot \frac{2x}{x^2 + y^2 + z^2} = \frac{x}{x^2 + y^2 + z^2}
\]
By symmetry:
\[
\frac{\partial u}{\partial y} = \frac{y}{x^2 + y^2 + z^2}, \quad \frac{\partial u}{\partial z} = \frac{z}{x^2 + y^2 + z^2}
\]

Step 4: Calculate the second partial derivatives.
Differentiating $\frac{\partial u}{\partial x}$ with respect to $x$ using the quotient rule:
\[
\frac{\partial^2 u}{\partial x^2} = \frac{(x^2 + y^2 + z^2)(1) - x(2x)}{(x^2 + y^2 + z^2)^2} = \frac{y^2 + z^2 - x^2}{(x^2 + y^2 + z^2)^2}
\]
By symmetry:
\[
\frac{\partial^2 u}{\partial y^2} = \frac{x^2 + z^2 - y^2}{(x^2 + y^2 + z^2)^2}, \quad \frac{\partial^2 u}{\partial z^2} = \frac{x^2 + y^2 - z^2}{(x^2 + y^2 + z^2)^2}
\]

Step 5: Sum the second partial derivatives.
\[
\frac{\partial^2 u}{\partial x^2} + \frac{\partial^2 u}{\partial y^2} + \frac{\partial^2 u}{\partial z^2} = \frac{(y^2 + z^2 - x^2) + (x^2 + z^2 - y^2) + (x^2 + y^2 - z^2)}{(x^2 + y^2 + z^2)^2}
\]
\[
= \frac{x^2 + y^2 + z^2}{(x^2 + y^2 + z^2)^2} = \frac{1}{x^2 + y^2 + z^2}
\]

Step 6: Multiply both sides by $(x^2 + y^2 + z^2)$.
\[
(x^2 + y^2 + z^2)\left(\frac{\partial^2 u}{\partial x^2} + \frac{\partial^2 u}{\partial y^2} + \frac{\partial^2 u}{\partial z^2}\right) = (x^2 + y^2 + z^2) \cdot \frac{1}{x^2 + y^2 + z^2} = 1
\]
Hence proved.

Final Answer:
\[
\boxed{(x^2 + y^2 + z^2)\left(\frac{\partial^2 u}{\partial x^2} + \frac{\partial^2 u}{\partial y^2} + \frac{\partial^2 u}{\partial z^2}\right) = 1}
\]
%%%%%%%%%%%%%%%%%%%%%%%%%%%%%%%%%%%%%%%%%%%%%%%%%%
% QUESTION_END
%%%%%%%%%%%%%%%%%%%%%%%%%%%%%%%%%%%%%%%%%%%%%%%%%%

\end{document}